\begin{table}[H]
\fontsize{6.0pt}{7.2pt}\selectfont
\begin{tabular*}{\linewidth}{@{\extracolsep{\fill}}lcccc}
\toprule
Region & Fire Size km²,
mean (SD) & WUI
(\% in region) & Fatalities per 100K
(total) & Population Affected Annually, mean (SD) \\ 
\midrule\addlinespace[2.5pt]
Nationwide & 44.9 (189.6) & 3,652 (61.8\%) &    0.159 & 4,921,855 (2,360,064) \\ 
Alaska & 144.3 (384.0) &     1 (2.8\%) &    0.000 & 3,300 (4,234) \\ 
California & 94.6 (320.6) &   584 (71.6\%) &    0.577 & 2,174,148 (1,747,526) \\ 
Eastern & 2.8 (21.0) &   371 (57.4\%) &    0.034 & 385,942 (339,106) \\ 
Hawaii & 24.2 (38.8) &    15 (83.3\%) &    6.734 & 98,789 (103,394) \\ 
Intermountain & 109.4 (253.5) &   155 (47.0\%) &    0.240 & 215,154 (226,170) \\ 
Northern Rockies & 92.2 (187.2) &    91 (40.6\%) &    0.179 & 30,756 (31,725) \\ 
Pacific Northwest & 129.0 (251.4) &   217 (63.6\%) &    0.078 & 141,326 (206,336) \\ 
Southern & 12.4 (109.7) & 1,948 (64.6\%) &    0.066 & 1,678,798 (1,282,424) \\ 
Southern Rockies & 61.9 (128.0) &   145 (52.9\%) &    0.216 & 174,297 (335,240) \\ 
Southwestern & 119.4 (310.0) &   125 (59.0\%) &    0.227 & 111,596 (141,707) \\ 
\bottomrule
\end{tabular*}
\caption{Counts of wildfire burn zone disasters, structures destroyed, and civilian fatalities by modified US Forest Service region, 2000--2024. Percentages represent each region's share of the national total.}
\label{tab:fire_description}
\end{table}

